% =====================================================================
%  ALANYA — Dossier de conception
%  Étape 4 : comptes business
%
%  Compilation :  latexmk -pdf compte-business.tex
% =====================================================================
\documentclass[11pt,a4paper,oneside]{report}

\usepackage{alanya-conception}
\setAlanyaDocTitre{Comptes business}

\hypersetup{
  colorlinks=true,
  linkcolor=AlanyaIndigoStrong,
  urlcolor=AlanyaIndigoStrong,
  pdftitle={Alanya - Comptes business},
  pdfauthor={Equipe Alanya},
  pdfsubject={Conception - profil business, catalogue, annuaire, tableau de bord}
}

\begin{document}

% =====================================================================
\begin{titlepage}
\thispagestyle{empty}
\begin{tikzpicture}[remember picture, overlay]
  \fill[AlanyaIndigo]
    (current page.north west) rectangle ([yshift=-6.4cm]current page.north east);
  \fill[AlanyaIndigoDark]
    ([yshift=-6.4cm]current page.north west) rectangle ([yshift=-6.0cm]current page.north east);
\end{tikzpicture}

\vspace*{1.1cm}
{\color{white}\sffamily
\begin{center}
  {\small\bfseries DOSSIER DE CONCEPTION}\\[6pt]
  {\footnotesize\color{AlanyaContainer} Document de travail --- à valider avant développement}
\end{center}}

\vspace{3.4cm}

\begin{center}\sffamily
  {\small\bfseries ÉTAPE 4}\\[10pt]
  {\color{AlanyaLine}\rule{0.6\linewidth}{0.6pt}}\\[20pt]
  {\fontsize{32}{36}\selectfont\bfseries\color{AlanyaIndigoDark} Comptes business}\\[10pt]
  {\large\color{AlanyaGray} Une vitrine, pas une boutique}\\[6pt]
  {\normalsize\itshape\color{AlanyaGray}
   Fiche étendue, catalogue, horaires, annuaire de proximité}\\[22pt]
  {\color{AlanyaLine}\rule{0.6\linewidth}{0.6pt}}
\end{center}

\vspace{1.2cm}

\begin{center}
  \badgePanier[3.2]\hspace{1.6cm}\badgePastillePanier[3.2]
\end{center}

\vfill

\begin{center}\sffamily
  \begin{tabular}{@{}r@{\quad}l@{}}
    \thd{Périmètre}   & Talky (Flutter) \textbullet{} Alanya-Backend \textbullet{} Alanya Admin\\[3pt]
    \thd{Statut}      & Conception --- aucune ligne de code écrite\\[3pt]
    \thd{Dépend de}   & Socle de compte (migration 023)\\[3pt]
    \thd{Hors périmètre} & Encaissement dans l'application\\[3pt]
    \thd{Date}        & \today\\
  \end{tabular}
\end{center}

\vspace{0.8cm}
\begin{center}\color{AlanyaGray}\small\sffamily
  Document confidentiel --- diffusion restreinte
\end{center}
\end{titlepage}

% =====================================================================
\chapter*{Avant-propos}
\addcontentsline{toc}{chapter}{Avant-propos}
\markboth{Avant-propos}{}
\thispagestyle{plain}

Le socle de compte (étape~1) a introduit le genre \emph{business} et son badge. Ce document
décrit ce que ce genre de compte \textbf{contient}~: sa fiche, son catalogue, ses horaires,
sa place dans un annuaire, et le tableau de bord qui lui rend compte de son audience.

\begin{note}[Une décision de périmètre à garder en tête]
Alanya reste une \textbf{vitrine}. Le commerçant présente son activité~; la transaction se
fait hors application. Les \guillemotleft~moyens de paiement~\guillemotright{} de sa fiche
sont une information affichée, pas un encaissement.
\end{note}

\tableofcontents

% =====================================================================
\part{Approche naïve}
% =====================================================================

\chapter{Ce qu'un commerçant peut présenter}

\section{La fiche}

Un compte business dispose d'une fiche étendue, absente d'un compte personnel~:

\begin{center}
\begin{tabularx}{\linewidth}{@{}p{4.4cm}X@{}}
\toprule
\thd{Identité} & Raison sociale, description, logo, image de couverture. \\
\thd{Catégorie} & Une catégorie principale, choisie dans une liste courte. \\
\thd{Mots-clés} & Saisis librement par le commerçant, pour être trouvé. \\
\thd{Lieu} & Adresse, position sur la carte, distance affichée au visiteur. \\
\thd{Horaires} & Jours et créneaux d'ouverture, exceptions. \\
\thd{Liens} & Site web, réseaux sociaux, dans l'ordre choisi. \\
\thd{Médias} & Images et vidéos de présentation. \\
\thd{Moyens de paiement} & Ce que l'établissement accepte sur place. \\
\thd{Catalogue} & Produits ou services, avec prix et disponibilité. \\
\bottomrule
\end{tabularx}
\end{center}

Trois de ces rubriques méritent d'être détaillées, parce qu'elles ne sont pas de simples
champs de formulaire.

\section{Le lieu et la géolocalisation}

C'est la rubrique dont dépend tout l'annuaire, et la seule qui parte réellement de zéro ---
Alanya ne stocke aujourd'hui aucune donnée géographique.

\begin{center}
\begin{tabularx}{\linewidth}{@{}p{4.4cm}X@{}}
\toprule
\thd{Adresse écrite} & Ce que le visiteur lit et recopie.
\guillemotleft~Avenue Kennedy, Yaoundé~\guillemotright. Texte libre~: il n'existe pas de
référentiel d'adresses fiable sur le marché visé. \\
\thd{Position} & Un point sur la carte, posé par le commerçant lui-même. C'est
\emph{elle} qui sert à la recherche par proximité, pas l'adresse écrite. \\
\thd{Distance affichée} & Calculée à la volée entre le visiteur et le commerce.
\guillemotleft~à 1,2~km de vous~\guillemotright. \\
\bottomrule
\end{tabularx}
\end{center}

\begin{note}[Comment le commerçant pose son point]
Deux gestes seulement~: \guillemotleft~utiliser ma position actuelle~\guillemotright{}
quand il est sur place, ou déplacer un repère sur une carte. Pas de géocodage d'adresse ---
il donnerait des résultats médiocres et masquerait les erreurs derrière une fausse
précision. Le commerçant sait où est sa boutique~; on lui demande de le montrer.
\end{note}

\begin{constat}[Aucune dépendance à ajouter]
\code{geolocator} et \code{flutter\_map} sont \textbf{déjà} dans \code{pubspec.yaml} --- ils
servent au partage de position dans les conversations
(\code{lib/screens/chats/location\_picker\_screen.dart}). Le sélecteur de position du
commerçant peut réutiliser cet écran.
\end{constat}

\section{Médias de présentation}

Images \emph{et} vidéos. C'est ce qui fait la différence entre une fiche remplie et une
fiche convaincante --- et c'est aussi la rubrique la plus coûteuse à exploiter.

\begin{center}
\begin{tabularx}{\linewidth}{@{}p{4.4cm}X@{}}
\toprule
\thd{Couverture} & Une image en tête de fiche, format paysage. \\
\thd{Galerie} & Photos de l'établissement, des produits, de l'équipe. Ordre choisi par le
commerçant. \\
\thd{Vidéos} & Présentation de l'activité. Avec vignette, durée, et lecture au clic. \\
\thd{Logo} & Réutilise l'avatar du compte --- pas de champ supplémentaire. \\
\bottomrule
\end{tabularx}
\end{center}

\begin{alerte}[La vidéo est un chantier, pas un champ]
Une vidéo téléversée depuis un téléphone pèse couramment plusieurs dizaines de mégaoctets.
La servir telle quelle à chaque ouverture de fiche est intenable. Il faut une vignette, une
limite de durée, et un transcodage --- lequel suppose la file de jobs conçue à l'étape~3.
Voir le chapitre~9.
\end{alerte}

\section{Liens et présence en ligne}

Site web, Facebook, Instagram, WhatsApp, TikTok\ldots{} Le commerçant en ajoute autant qu'il
veut et choisit leur ordre d'affichage.

\begin{note}[Pas une colonne par réseau]
Le point que le document d'amorce soulignait à juste titre. Une table
\code{business\_link (type, url, ordre)} permet d'ajouter le réseau à la mode de l'an
prochain sans migration. Une colonne \code{facebook\_url} ne le permettrait pas.
\end{note}

\section{Les quatre écrans}

\begin{center}
  \imageOuCadre{maquettes/comptes-business-ecrans.png}{\linewidth}%
    {Fiche business, catalogue, éditeur d'horaires, annuaire\\
     \texttt{maquettes/comptes-business-ecrans.png}}
\end{center}

\begin{note}[Maquette interactive]
\href{https://claude.ai/code/artifact/37b8e6c7-1235-42df-aabc-d66b8f26ff67}%
{\texttt{claude.ai/code/artifact/37b8e6c7}}\\
Source versionnée~: \code{docs/architecture/maquettes/comptes-business.html}
\end{note}

\section{Le badge fait son travail}

Dans l'annuaire, trois commerces se distinguent d'un coup d'{\oe}il~: deux portent la
pastille pleine \badgePastillePanier[1] --- vérifiés par Alanya --- et un porte le panier nu
\badgePanier[1], c'est-à-dire un commerce qui s'est déclaré tel sans qu'on l'ait contrôlé.

C'est exactement la distinction que le socle devait rendre lisible à 13~pixels, et c'est
dans l'annuaire qu'elle compte le plus~: un visiteur qui cherche une pharmacie de garde
prend une décision sur cette base.

\chapter{La catégorie et les mots-clés}

\section{Une liste courte, pas une arborescence}

Le choix retenu~: \textbf{une catégorie principale} dans une liste courte, plus des
\textbf{mots-clés libres} saisis par le commerçant.

\begin{center}
\begin{tabularx}{\linewidth}{@{}p{4.4cm}X@{}}
\toprule
\thd{Pourquoi pas deux niveaux} & Un sélecteur à deux étages est pénible sur un téléphone,
et il faudrait arrêter dès aujourd'hui une liste longue qu'on ne pourra plus changer. \\
\thd{Pourquoi des mots-clés} & Ils absorbent tout ce que la liste ne prévoit pas ---
\guillemotleft~garde de nuit~\guillemotright, \guillemotleft~livraison~\guillemotright,
\guillemotleft~sur mesure~\guillemotright{} --- sans figer quoi que ce soit. \\
\bottomrule
\end{tabularx}
\end{center}

\begin{alerte}[Le revers, à assumer]
La qualité de la recherche dépendra de ce que les commerçants écrivent. Deux garde-fous
tempèrent~: les mots-clés sont \emph{normalisés} (minuscules, accents repliés) et le champ
propose les mots déjà utilisés par d'autres commerces de la même catégorie. Un mot-clé
partagé vaut mieux que dix orthographes du même mot.
\end{alerte}

\section{Une taxonomie est difficile à changer}

C'est le point que le document d'amorce soulignait à juste titre~: une fois des milliers de
fiches rangées, renommer ou fusionner une catégorie devient une opération de reprise de
données. D'où la liste courte --- vingt à trente entrées --- et la table de référence
éditable plutôt que des valeurs figées dans le code.

\chapter{Les horaires, l'écran le plus difficile}

Ce n'est pas un champ de formulaire. Six cas doivent tenir sur un écran de téléphone, et
cinq d'entre eux restent invisibles tant qu'on n'a pas essayé.

\begin{center}
\begin{tabularx}{\linewidth}{@{}p{4.1cm}X@{}}
\toprule
\thd{Multi-créneaux} & Ouvert le matin, fermé à midi, rouvert l'après-midi. Deux plages pour
un seul jour --- le cas le plus courant, et celui qu'un champ \guillemotleft~de\ldots{}
à\ldots{}~\guillemotright{} ne sait pas exprimer. \\
\thd{24\,h/24} & N'est pas \guillemotleft~00:00 $\rightarrow$ 00:00~\guillemotright. C'est un
état distinct, sans quoi le calcul \guillemotleft~ouvert maintenant~\guillemotright{} se
trompe. \\
\thd{Passage de minuit} & Un bar ouvert de 20\,h à 2\,h~: l'heure de fin est
\emph{avant} l'heure de début. \\
\thd{Copie d'un jour} & Saisir sept fois les mêmes horaires fait abandonner.
\guillemotleft~Copier vers~\guillemotright{} n'est pas un confort, c'est ce qui décide si la
fiche sera remplie. \\
\thd{Exceptions} & Jours fériés, congés, fermeture ponctuelle. Une liste à part, pas un
champ de plus. \\
\thd{Fuseau horaire} & \guillemotleft~Ouvert maintenant~\guillemotright{} se calcule dans le
fuseau du commerce, pas dans celui du visiteur. \\
\bottomrule
\end{tabularx}
\end{center}

\begin{constat}[Le fuseau existe déjà]
La table \code{pays} porte \code{timeZone} et \code{decalageHoraire}
(\code{migrations/001\_initial\_schema.sql}), et \code{users.idPays} les relie. Rien à
créer --- mais tout à ne pas oublier.
\end{constat}

\begin{note}[À budgéter comme un écran]
Cet éditeur mérite sa propre estimation, au même titre qu'un écran de conversation. Le
traiter comme un champ du formulaire de profil garantit qu'il sera bâclé, puis refait.
\end{note}

\chapter{Le tableau de bord du commerçant}

\section{Ce qu'il montre}

Puisque les transactions se font hors application, le tableau de bord ne peut pas parler
d'argent. Il parle de ce qu'Alanya \textbf{observe réellement}~: l'audience de la fiche.

\begin{itemize}[leftmargin=1.4em]
  \item nombre de vues du profil, et son évolution~;
  \item provenance des visites --- annuaire, recherche, conversation, lien partagé~;
  \item répartition géographique des visiteurs~;
  \item conversations entamées depuis la fiche.
\end{itemize}

\begin{note}[Pourquoi pas de commandes déclarées]
On aurait pu laisser le commerçant saisir ses ventes pour obtenir un
\guillemotleft~suivi des transactions~\guillemotright. L'expérience de ce genre de saisie
manuelle est constante~: elle est abandonnée après quelques semaines, et le tableau de bord
devient faux --- pire que vide. Écarté volontairement.
\end{note}

\section{Ce que cela implique}

Compter des vues suppose de les enregistrer. C'est un journal d'événements, avec le volume
que cela suppose --- voir le chapitre technique~8.

% =====================================================================
\part{Approche technique}
% =====================================================================

\chapter{Modèle de données}

\section{Vue d'ensemble}

\begin{center}
\begin{tikzpicture}[
  x=1cm, y=1cm,
  t/.style={draw=AlanyaLine, fill=white, rounded corners=3pt,
            font=\sffamily\scriptsize, inner sep=5pt, text=AlanyaInk, align=center},
  socle/.style={t, draw=AlanyaIndigo, line width=1pt, fill=AlanyaContainer,
                font=\sffamily\small\bfseries, text=AlanyaIndigoDark},
  ref/.style={t, draw=AlanyaMute, fill=AlanyaSurface, text=AlanyaGray},
  l/.style={draw=AlanyaLine, semithick}
]
  \node[socle] (u)  at (0,0)      {users};
  \node[t]     (bp) at (0,-2.0)   {business\_profile};
  \node[t]     (h)  at (-4.3,-3.9){business\_hours};
  \node[t]     (ex) at (-4.3,-5.2){business\_hours\_exception};
  \node[t]     (li) at (0,-3.9)   {business\_link};
  \node[t]     (me) at (0,-5.2)   {business\_media};
  \node[t]     (pr) at (4.3,-3.9) {business\_product};
  \node[t]     (pv) at (4.3,-5.2) {business\_profile\_view};
  \node[ref]   (ca) at (5.2,-2.0) {business\_category\\[-1pt]{\itshape référence}};
  \node[ref]   (kw) at (-5.0,-2.0){business\_keyword};

  \draw[l] (u)  -- (bp);
  \draw[l] (bp) -- (ca);
  \draw[l] (bp) -- (kw);
  \foreach \n in {h,li,me,pr,pv} { \draw[l] (bp) -- (\n); }
  \draw[l] (h) -- (ex);
\end{tikzpicture}
\end{center}

\section{Le profil}

\begin{lstlisting}[language=SQL, caption={\code{migrations/026\_business.sql} --- extrait}]
-- 026_business.sql -- profil business etendu
-- Requiert 023 (socle de compte).

CREATE TABLE IF NOT EXISTS business_category (
  id       SMALLINT     NOT NULL AUTO_INCREMENT,
  slug     VARCHAR(40)  NOT NULL,
  libelle  VARCHAR(80)  NOT NULL,
  ordre    SMALLINT     NOT NULL DEFAULT 0,
  actif    TINYINT      NOT NULL DEFAULT 1,
  PRIMARY KEY (id),
  UNIQUE KEY uq_cat_slug (slug)
) ENGINE=InnoDB DEFAULT CHARSET=utf8mb4 COLLATE=utf8mb4_unicode_ci;

CREATE TABLE IF NOT EXISTS business_profile (
  alanyaID     INT          NOT NULL,        -- 1-1 avec le compte
  legal_name   VARCHAR(120) NOT NULL,
  description  TEXT         NULL,
  category_id  SMALLINT     NULL,
  cover_url    VARCHAR(255) NULL,
  address      VARCHAR(255) NULL,
  -- SPATIAL INDEX impose NOT NULL : on prevoit la colonne des maintenant,
  -- l'ajouter apres coup sur une table remplie est penible.
  geo          POINT        NOT NULL SRID 4326,
  timezone     VARCHAR(64)  NULL,            -- sinon deduit via users.idPays
  is_hidden    TINYINT      NOT NULL DEFAULT 0
    COMMENT 'Profil conserve mais masque (retour en compte personnel)',
  created_at   DATETIME     NOT NULL DEFAULT CURRENT_TIMESTAMP,
  updated_at   DATETIME     NOT NULL DEFAULT CURRENT_TIMESTAMP
                            ON UPDATE CURRENT_TIMESTAMP,
  PRIMARY KEY (alanyaID),
  SPATIAL INDEX idx_bp_geo (geo),
  KEY idx_bp_categorie (category_id, is_hidden),
  CONSTRAINT fk_bp_user FOREIGN KEY (alanyaID)
    REFERENCES users(alanyaID) ON DELETE CASCADE,
  CONSTRAINT fk_bp_cat  FOREIGN KEY (category_id)
    REFERENCES business_category(id)
) ENGINE=InnoDB DEFAULT CHARSET=utf8mb4 COLLATE=utf8mb4_unicode_ci;
\end{lstlisting}

\begin{alerte}[Deux contraintes MySQL à ne pas découvrir en route]
\code{SPATIAL INDEX} exige que la colonne soit \code{NOT NULL} et porte un SRID explicite.
Un commerce sans coordonnées connues doit donc recevoir une position par défaut --- celle du
chef-lieu de son pays --- plutôt qu'un \code{NULL}. Et \code{utf8mb4} est indispensable~: le
premier emoji dans un nom d'entreprise ferait échouer l'insertion sur un jeu plus étroit.
\end{alerte}

\section{Liens, médias, mots-clés}

\begin{lstlisting}[language=SQL, caption={Tables satellites}]
-- Pas une colonne par reseau social : ajouter TikTok demain ne doit
-- pas demander de migration.
CREATE TABLE IF NOT EXISTS business_link (
  id       BIGINT       NOT NULL AUTO_INCREMENT,
  alanyaID INT          NOT NULL,
  kind     TINYINT      NOT NULL COMMENT '0=site 1=facebook 2=instagram ...',
  url      VARCHAR(255) NOT NULL,
  ordre    SMALLINT     NOT NULL DEFAULT 0,
  PRIMARY KEY (id),
  KEY idx_bl_owner (alanyaID, ordre),
  CONSTRAINT fk_bl_user FOREIGN KEY (alanyaID)
    REFERENCES business_profile(alanyaID) ON DELETE CASCADE
) ENGINE=InnoDB DEFAULT CHARSET=utf8mb4;

CREATE TABLE IF NOT EXISTS business_keyword (
  alanyaID   INT         NOT NULL,
  keyword    VARCHAR(40) NOT NULL COMMENT 'Deja normalise : minuscules, accents replies',
  PRIMARY KEY (alanyaID, keyword),
  KEY idx_bk_mot (keyword),
  CONSTRAINT fk_bk_user FOREIGN KEY (alanyaID)
    REFERENCES business_profile(alanyaID) ON DELETE CASCADE
) ENGINE=InnoDB DEFAULT CHARSET=utf8mb4;

-- Images ET videos de presentation. Une seule table : la difference tient
-- au champ `kind` et aux colonnes de transcodage, nulles pour une image.
CREATE TABLE IF NOT EXISTS business_media (
  id           BIGINT       NOT NULL AUTO_INCREMENT,
  alanyaID     INT          NOT NULL,
  kind         TINYINT      NOT NULL DEFAULT 0 COMMENT '0=image 1=video',
  role         TINYINT      NOT NULL DEFAULT 0 COMMENT '0=galerie 1=couverture',
  url          VARCHAR(255) NOT NULL,
  thumb_url    VARCHAR(255) NULL  COMMENT 'Vignette : obligatoire pour une video',
  width        SMALLINT     NULL,
  height       SMALLINT     NULL,
  duration_ms  INT          NULL  COMMENT 'Video uniquement',
  size_bytes   BIGINT       NULL,
  -- Cycle de vie du transcodage, pilote par la file de jobs (etape 3).
  status       TINYINT      NOT NULL DEFAULT 0
                 COMMENT '0=en attente 1=pret 2=echec',
  ordre        SMALLINT     NOT NULL DEFAULT 0,
  created_at   DATETIME     NOT NULL DEFAULT CURRENT_TIMESTAMP,
  PRIMARY KEY (id),
  KEY idx_bm_owner (alanyaID, role, ordre),
  KEY idx_bm_statut (status),
  CONSTRAINT fk_bm_user FOREIGN KEY (alanyaID)
    REFERENCES business_profile(alanyaID) ON DELETE CASCADE
) ENGINE=InnoDB DEFAULT CHARSET=utf8mb4;
\end{lstlisting}

\begin{note}[Une seule table pour les images et les vidéos]
Les séparer obligerait à fusionner deux listes à chaque affichage de galerie, et à dupliquer
l'ordre d'affichage. \code{kind} distingue les deux~; les colonnes propres à la vidéo
(\code{duration\_ms}, \code{thumb\_url}, \code{status}) restent nulles pour une image.
\end{note}

\begin{note}[\code{status} n'est pas décoratif]
Une vidéo est \emph{visible mais pas encore lisible} tant que le transcodage n'a pas
abouti. Sans cet état, la fiche afficherait une vidéo qui ne démarre pas --- et le
commerçant conclurait que l'application est cassée.
\end{note}

\section{Horaires}

\begin{lstlisting}[language=SQL, caption={Horaires et exceptions}]
CREATE TABLE IF NOT EXISTS business_hours (
  id         BIGINT  NOT NULL AUTO_INCREMENT,
  alanyaID   INT     NOT NULL,
  weekday    TINYINT NOT NULL COMMENT '1=lundi ... 7=dimanche',
  -- Minutes depuis minuit : 480 = 08:00. Un TIME suffirait, mais l'entier
  -- rend le passage de minuit trivial (fin < debut => le lendemain).
  start_min  SMALLINT NULL,
  end_min    SMALLINT NULL,
  all_day    TINYINT NOT NULL DEFAULT 0 COMMENT '1 = 24h/24, start/end ignores',
  PRIMARY KEY (id),
  KEY idx_bh_jour (alanyaID, weekday),
  CONSTRAINT fk_bh_user FOREIGN KEY (alanyaID)
    REFERENCES business_profile(alanyaID) ON DELETE CASCADE
) ENGINE=InnoDB DEFAULT CHARSET=utf8mb4;

-- Absence de ligne pour un jour = ferme ce jour-la.
-- Plusieurs lignes pour un meme jour = plusieurs creneaux.

CREATE TABLE IF NOT EXISTS business_hours_exception (
  id        BIGINT      NOT NULL AUTO_INCREMENT,
  alanyaID  INT         NOT NULL,
  day       DATE        NOT NULL,
  closed    TINYINT     NOT NULL DEFAULT 1,
  start_min SMALLINT    NULL,   -- horaires speciaux si closed = 0
  end_min   SMALLINT    NULL,
  label     VARCHAR(80) NULL,   -- "fete nationale"
  PRIMARY KEY (id),
  UNIQUE KEY uq_bhe (alanyaID, day),
  CONSTRAINT fk_bhe_user FOREIGN KEY (alanyaID)
    REFERENCES business_profile(alanyaID) ON DELETE CASCADE
) ENGINE=InnoDB DEFAULT CHARSET=utf8mb4;
\end{lstlisting}

\begin{note}[Trois conventions qui évitent des cas particuliers]
\textbf{Pas de ligne} pour un jour signifie \emph{fermé} --- inutile de stocker l'absence.
\textbf{Plusieurs lignes} pour le même jour donnent les créneaux multiples, sans colonne
supplémentaire. Et \code{end\_min < start\_min} signifie \emph{le lendemain}, ce qui règle
le passage de minuit sans champ dédié.
\end{note}

\section{Catalogue}

\begin{lstlisting}[language=SQL, caption={Produits}]
CREATE TABLE IF NOT EXISTS business_product (
  id           BIGINT       NOT NULL AUTO_INCREMENT,
  alanyaID     INT          NOT NULL,
  name         VARCHAR(120) NOT NULL,
  description  TEXT         NULL,
  -- Entier en unite mineure + code devise. JAMAIS un flottant.
  price_minor  BIGINT       NULL,
  currency     CHAR(3)      NOT NULL DEFAULT 'XAF',
  unit         VARCHAR(40)  NULL COMMENT 'boite de 20, 500 ml',
  media_url    VARCHAR(255) NULL,
  availability TINYINT      NOT NULL DEFAULT 0
                 COMMENT '0=disponible 1=sur commande 2=epuise',
  ordre        SMALLINT     NOT NULL DEFAULT 0,
  is_hidden    TINYINT      NOT NULL DEFAULT 0,
  created_at   DATETIME     NOT NULL DEFAULT CURRENT_TIMESTAMP,
  PRIMARY KEY (id),
  KEY idx_bpr_owner (alanyaID, is_hidden, ordre),
  CONSTRAINT fk_bpr_user FOREIGN KEY (alanyaID)
    REFERENCES business_profile(alanyaID) ON DELETE CASCADE
) ENGINE=InnoDB DEFAULT CHARSET=utf8mb4 COLLATE=utf8mb4_unicode_ci;
\end{lstlisting}

\begin{alerte}[Le prix n'est jamais un flottant]
1\,200~FCFA se stocke \code{120000} en unité mineure, avec \code{XAF} à côté. Un
\code{FLOAT} ou un \code{DOUBLE} sur des prix finit toujours par produire un centime
fantôme, et le franc CFA n'a pas de subdivision --- raison de plus pour que l'unité soit
explicite plutôt que devinée.
\end{alerte}

\chapter{Un lieu, ou plusieurs}

\section{Ce qui est retenu}

\textbf{Un seul lieu par compte.} L'adresse, la position et les horaires sont portés par
\code{business\_profile}. Une chaîne de trois boutiques crée aujourd'hui trois comptes.

\section{La variante multi-lieux}

Elle est documentée ici parce qu'elle est envisagée comme \textbf{fonctionnalité payante} et
qu'il vaut mieux en connaître le coût avant de la promettre.

\begin{center}
\begin{tabularx}{\linewidth}{@{}p{4.2cm}XX@{}}
\toprule
 & \thd{Un lieu (retenu)} & \thd{Plusieurs lieux} \\
\midrule
Adresse, position, horaires & sur \code{business\_profile} &
  déplacés dans \code{business\_location} \\
Annuaire & cherche des \emph{comptes} & cherche des \emph{lieux}, remonte le compte \\
Index spatial & sur \code{business\_profile} & sur \code{business\_location} \\
Éditeur d'horaires & un jeu & un jeu \textbf{par lieu} \\
Fiche & une adresse & sélecteur de lieu en tête de fiche \\
\bottomrule
\end{tabularx}
\end{center}

\textbf{Coût de bascule}~: création de \code{business\_location}, migration des colonnes
d'adresse, de position et d'horaires depuis \code{business\_profile}, déplacement de l'index
spatial, et réécriture de la requête d'annuaire. Le reste --- catalogue, liens, médias,
mots-clés --- n'est pas touché.

\begin{note}[Comment se garder la porte ouverte à moindre frais]
Concevoir dès maintenant les requêtes d'annuaire pour renvoyer un couple
\emph{(compte, lieu)} plutôt qu'un compte seul. Avec un seul lieu, le second terme est
constant~; le jour de la bascule, seule la source change, pas les appelants.
\end{note}

\begin{alerte}[Le déblocage vient de l'abonnement, pas du badge]
Si les lieux multiples étaient conditionnés à \code{verification\_status = 2}, le badge
deviendrait une fonctionnalité qui se vend --- exactement le recouplage entre identité et
argent que le découplage de l'étape~5 cherche à éviter. La condition doit porter sur
\code{subscriptions}. En pratique l'utilisateur voit la même chose~; dans le modèle, la porte
n'est pas ouverte par la même clé.
\end{alerte}

\chapter{Annuaire et recherche de proximité}

\section{L'état actuel}

\begin{constat}[Il n'y a aucune donnée géographique]
Le schéma ne contient pas une seule colonne de position. La page
\guillemotleft~Géolocalisation~\guillemotright{} de l'administration est une table de
\textbf{54~pays codée en dur} (\code{talky-admin/lib/geo-utils.ts}), alimentée par des
\emph{noms de pays en chaîne} issus de \code{/api/admin/stats}. Il n'existe pas d'endpoint
géographique côté backend.
\end{constat}

L'annuaire de proximité part donc de zéro. Ce n'est pas une évolution de l'existant, c'est
une brique nouvelle.

\section{La requête}

\begin{lstlisting}[language=SQL, caption={Commerces ouverts, par proximité}]
SELECT b.alanyaID, u.nom, b.legal_name, c.libelle,
       ST_Distance_Sphere(b.geo, ST_SRID(POINT(?, ?), 4326)) AS metres
FROM business_profile b
JOIN users u ON u.alanyaID = b.alanyaID
LEFT JOIN business_category c ON c.id = b.category_id
WHERE b.is_hidden = 0
  AND u.exclus = 0
  AND MBRContains(
        ST_Buffer(ST_SRID(POINT(?, ?), 4326), ?),   -- prefiltre sur l'index
        b.geo)
  AND (? IS NULL OR b.category_id = ?)
ORDER BY metres
LIMIT 50;
\end{lstlisting}

\begin{note}[Pourquoi le préfiltre]
\code{ST\_Distance\_Sphere} est une fonction~: seule, elle impose un balayage complet.
L'index spatial ne sert que si la requête contient une contrainte de \emph{contenance}
(\code{MBRContains}) qu'il sait exploiter. On restreint donc d'abord à un rectangle, puis on
trie sur la distance exacte à l'intérieur.
\end{note}

\section{\guillemotleft~Ouvert maintenant~\guillemotright}

Ce filtre ne se calcule pas en SQL sans douleur~: il dépend du jour courant \emph{dans le
fuseau du commerce}, des créneaux multiples, du cas 24\,h/24, du passage de minuit et des
exceptions. On charge les horaires des résultats retenus --- au plus cinquante --- et on
évalue en mémoire. La même fonction sert à l'affichage
\guillemotleft~ferme à 19\,h~\guillemotright{} sur la fiche.

\chapter{Tableau de bord d'audience}

\section{Journal d'événements}

\begin{lstlisting}[language=SQL, caption={Vues de fiche}]
CREATE TABLE IF NOT EXISTS business_profile_view (
  id         BIGINT   NOT NULL AUTO_INCREMENT,
  alanyaID   INT      NOT NULL,          -- le commerce vu
  viewer_id  INT      NULL,              -- NULL si non connecte
  source     TINYINT  NOT NULL DEFAULT 0
               COMMENT '0=annuaire 1=recherche 2=conversation 3=lien partage',
  viewed_at  DATETIME NOT NULL DEFAULT CURRENT_TIMESTAMP,
  PRIMARY KEY (id),
  KEY idx_bpv_owner_date (alanyaID, viewed_at),
  CONSTRAINT fk_bpv_biz FOREIGN KEY (alanyaID)
    REFERENCES business_profile(alanyaID) ON DELETE CASCADE
) ENGINE=InnoDB DEFAULT CHARSET=utf8mb4;
\end{lstlisting}

\begin{alerte}[Une table qui grossit vite]
Chaque ouverture de fiche est une ligne. Il faut décider dès maintenant~: anti-rebond côté
serveur (une vue par visiteur et par heure), agrégation quotidienne dans une table de
synthèse, et purge du détail au-delà de quelques mois. Sans cela, cette table dépassera
\code{message} en volume.
\end{alerte}

\section{Ce que le commerçant voit}

Vues sur la période, évolution, répartition par provenance et par pays, conversations
entamées depuis la fiche. Tout est dérivé de cette table et de \code{conversation}, sans
saisie du commerçant.

\chapter{Médias : stockage et transcodage}

\section{Ce que l'existant sait déjà faire}

\begin{constat}[Vérifié dans le code]
\code{src/middleware/upload.js} accepte déjà les images (jpeg, png, webp, gif) et les vidéos
(mp4, webm, quicktime, 3gpp), plafonne à \textbf{50~Mo} par fichier, range par type dans
\code{uploads/media/\{images,audio,video,files\}} et sert le tout en statique sur
\code{/uploads}. Côté mobile, \code{video\_thumbnail} et \code{image\_thumbnail\_service}
produisent déjà des vignettes pour les messages.
\end{constat}

Le téléversement n'est donc pas à construire. Ce qui manque, c'est tout ce qui vient après.

\section{Trois problèmes que la fiche business introduit}

\begin{center}
\begin{tabularx}{\linewidth}{@{}p{4.2cm}X@{}}
\toprule
\thd{Le volume} & Un message vidéo est vu une fois par un destinataire. Une vidéo de fiche
est servie à \emph{chaque} ouverture de la fiche, à tous les visiteurs. Le rapport n'est pas
du même ordre. \\
\thd{Le poids d'origine} & 50~Mo est acceptable pour un envoi ponctuel entre deux personnes.
C'est intenable pour un média public consulté depuis un réseau mobile. \\
\thd{La diversité des formats} & Un enregistrement issu d'un téléphone Android bas de gamme
et un autre d'un iPhone récent n'ont ni le même codec, ni la même orientation, ni la même
cadence. \\
\bottomrule
\end{tabularx}
\end{center}

\section{La chaîne de traitement}

\begin{center}
\begin{tikzpicture}[
  x=1cm, y=1cm,
  e/.style={draw=AlanyaLine, fill=AlanyaSurface, rounded corners=3pt,
            font=\sffamily\scriptsize, inner sep=5pt, minimum height=0.9cm,
            text width=2.4cm, text=AlanyaInk, align=center},
  f/.style={e, draw=AlanyaIndigo, line width=0.9pt, fill=AlanyaContainer,
            text=AlanyaIndigoDark},
  a/.style={-{Stealth[length=2mm]}, draw=AlanyaGray, semithick},
  l/.style={font=\sffamily\scriptsize, text=AlanyaGray, fill=white, inner sep=2pt}
]
  \node[f] (up) at (0,0)     {Téléversement\\[-2pt]{\scriptsize fichier d'origine}};
  \node[e] (jq) at (3.9,0)   {File de jobs\\[-2pt]{\scriptsize étape 3}};
  \node[e] (tr) at (7.8,0)   {Transcodage\\[-2pt]{\scriptsize + vignette}};
  \node[f] (ok) at (11.7,0)  {\code{status = 1}\\[-2pt]{\scriptsize lisible}};

  \draw[a] (up) -- node[l, above] {\code{status = 0}} (jq);
  \draw[a] (jq) -- node[l, above] {consomme} (tr);
  \draw[a] (tr) -- node[l, above] {publie} (ok);

  \node[l, text=AlanyaMute, align=center, text width=4cm] at (7.8,-1.5)
    {échec après $N$ tentatives\\ $\rightarrow$ \code{status = 2}, média retiré de la fiche};
  \draw[draw=AlanyaLine, dashed] (7.8,-0.5) -- (7.8,-1.05);
\end{tikzpicture}
\end{center}

Le transcodage est le \textbf{deuxième client} de la file de jobs, après le fan-out des
notifications --- exactement ce que prévoyait le §7 du document d'amorce. C'est un argument
de plus pour que la file soit construite à l'étape~3 et non plus tard.

\section{Limites à fixer}

Aucune n'est technique~: ce sont des décisions produit, à arrêter avant l'implémentation.

\begin{center}
\begin{tabularx}{\linewidth}{@{}p{5.2cm}X@{}}
\toprule
Images de galerie & 10 par fiche~? \\
Vidéos de présentation & 2 par fiche, 60~secondes chacune~? \\
Résolution servie & 1080p suffit pour un téléphone~; conserver l'original~? \\
Poids après transcodage & viser quelques mégaoctets, pas quelques dizaines \\
\bottomrule
\end{tabularx}
\end{center}

\begin{alerte}[Le disque local va montrer ses limites]
Les médias de fiches sont publics, consultés en boucle et conservés indéfiniment --- à la
différence des médias de conversation. C'est le premier usage qui justifiera un stockage
objet et un cache en frontal. Le disque local tiendra le temps de la mise en service, pas
au-delà.
\end{alerte}

\chapter{Impact fichier par fichier}

\section{Backend}

\begin{center}
\begin{tabularx}{\linewidth}{@{}p{6.2cm}p{1.6cm}X@{}}
\toprule
\thd{Fichier} & \thd{Nature} & \thd{Objet} \\
\midrule
\code{migrations/026\_business.sql} & nouveau & Profil, catégories, liens, mots-clés, horaires, produits, vues \\
\code{src/controllers/businessController.js} & nouveau & Fiche, liens, médias, mots-clés \\
\code{src/controllers/businessHoursController.js} & nouveau & Créneaux et exceptions \\
\code{src/controllers/businessProductController.js} & nouveau & Catalogue \\
\code{src/controllers/directoryController.js} & nouveau & Annuaire, proximité, filtres \\
\code{src/utils/openingHours.js} & nouveau & Évaluation « ouvert maintenant », fuseau compris \\
\code{src/routes/business.js} & nouveau & \code{/api/business}, \code{/api/directory} \\
\code{server.js} & modifié & Montage des routes \\
\bottomrule
\end{tabularx}
\end{center}

\section{Mobile}

\begin{center}
\begin{tabularx}{\linewidth}{@{}p{6.2cm}X@{}}
\toprule
\thd{Fichier} & \thd{Objet} \\
\midrule
\code{lib/screens/business/business\_profile\_screen.dart} \emph{(nouveau)} & Fiche vue par un visiteur \\
\code{lib/screens/business/business\_edit\_screen.dart} \emph{(nouveau)} & Édition par le propriétaire \\
\code{lib/screens/business/hours\_editor\_screen.dart} \emph{(nouveau)} & \textbf{À budgéter à part} \\
\code{lib/screens/business/catalog\_screen.dart} \emph{(nouveau)} & Grille et fiche produit \\
\code{lib/screens/business/directory\_screen.dart} \emph{(nouveau)} & Annuaire et proximité \\
\code{lib/screens/business/dashboard\_screen.dart} \emph{(nouveau)} & Audience de la fiche \\
\code{lib/core/db/app\_database.dart} & Cache local du profil et du catalogue (Drift v16 $\rightarrow$ v17) \\
\bottomrule
\end{tabularx}
\end{center}

\begin{note}[Réutiliser ce qui existe]
\code{geolocator} et \code{flutter\_map} sont \textbf{déjà} dans \code{pubspec.yaml} --- ils
servent au partage de position dans les conversations. L'annuaire n'ajoute aucune
dépendance. Idem pour \code{cached\_network\_image} sur les médias du catalogue.
\end{note}

\chapter{Points de vigilance et questions ouvertes}

\section{Chantiers d'infrastructure nommés}

\begin{enumerate}[leftmargin=1.6em]
  \item \textbf{L'index spatial} --- aucune donnée géographique n'existe aujourd'hui. Colonne
        \code{NOT NULL}, SRID explicite, et une position par défaut pour les commerces sans
        coordonnées.
  \item \textbf{Le journal de vues} --- anti-rebond, agrégation, purge. À concevoir avec la
        table, pas après.
  \item \textbf{Le transcodage vidéo} --- deuxième client de la file de jobs de l'étape~3.
        Sans lui, les vidéos de présentation sont inutilisables sur un réseau mobile.
  \item \textbf{Le stockage des médias} --- galerie, vidéos et catalogue multiplient les
        fichiers, publics et servis en boucle. Le disque local avec \code{multer}
        (\code{src/middleware/upload.js}) tiendra le temps de la mise en service~; c'est le
        premier candidat au stockage objet.
\end{enumerate}

\section{Questions ouvertes}

\begin{enumerate}[leftmargin=1.6em]
  \item La liste des 20 à 30 catégories reste à arrêter. Elle est difficile à changer
        ensuite~: elle mérite une revue dédiée.
  \item Un compte business a-t-il des réglages de confidentialité différents~? Sa fiche est
        publique par nature --- le \guillemotleft~vu à~\guillemotright{} et le
        \guillemotleft~en ligne~\guillemotright{} ont-ils encore un sens~?
  \item Quelles limites~: nombre de produits, de médias, de mots-clés~? Sans plafond, une
        fiche peut devenir une charge à elle seule.
  \item Un commerce peut-il être signalé~? Une arnaque commerciale n'est pas le même
        traitement qu'un signalement d'utilisateur, et rien n'existe aujourd'hui.
  \item Que devient le catalogue au retour en compte personnel~? Le champ
        \code{is\_hidden} le conserve --- confirmer que c'est bien le comportement voulu.
\end{enumerate}

\vfill
\begin{center}\small\color{AlanyaGray}\sffamily
Fin du dossier --- étape 4. \\
Étape suivante : certification.
\end{center}

\end{document}
